% ---------------------------------------------------------------------
%  Introduction générale (première page en chiffres arabes)
% ---------------------------------------------------------------------
\pagetitre{INTRODUCTION GÉNÉRALE}

Partir en voyage entre amis, organiser une sortie ou partager un appartement amène
toujours la même question : qui a payé quoi, et qui doit combien à qui ? Les dépenses sont
avancées par des personnes différentes, parfois dans plusieurs devises, et rarement
réparties à parts égales. Après quelques jours, les comptes deviennent flous, et le calcul
final finit souvent sur un coin de feuille ou dans un tableur partagé. Des applications
comme Tricount~\cite{tricount} ou Splitwise~\cite{splitwise} ont simplifié ce calcul, mais elles reposent toujours sur une
saisie manuelle de chaque dépense, et n'aident en rien à préparer le voyage lui-même.

Dans le même temps, les progrès des grands modèles de langage ouvrent de nouvelles
possibilités. Un ticket de caisse, une phrase comme « J'ai payé 45 € de taxi pour Sara et
moi » ou une liste de lieux à visiter contiennent déjà toute l'information utile : un
modèle bien entraîné peut l'extraire et la transformer en données structurées, tout en
restant assez compact pour être affiné sur un GPU grand public et exécuté sur un portable.

C'est dans ce contexte que s'inscrit ce projet, réalisé dans le cadre du module de
Génération Textuelle de la 5\textsuperscript{ème} année à l'EMSI. Il vise à concevoir et à
développer \equipay, une application mobile de partage des dépenses de groupe qui utilise
l'intelligence artificielle là où elle apporte un gain concret : un assistant
conversationnel pour accompagner l'utilisateur, un scanner de tickets, la saisie d'une
dépense en une phrase et l'organisation automatique du programme d'un voyage. L'application
repose sur Flutter pour le mobile, sur Supabase pour les comptes et les données partagées,
et sur un backend FastAPI qui sert les modèles en local. Ces modèles sont des versions
affinées de Qwen3, suivies dans MLflow depuis la préparation des données jusqu'au choix de
la version mise en production.

Ce rapport est organisé en cinq chapitres. Le premier présente le cadre du projet et la
démarche suivie, le deuxième l'étude préalable et la spécification des besoins, et le
troisième la conception de l'application. Le quatrième détaille la réalisation, en
particulier l'entraînement et l'évaluation des modèles, et le dernier présente
l'application à travers ses principaux écrans.
